\scanpage{47}
\mspage{45}
% archive sheet number: 46
l'action de $I$ sur $H^i_!(X_{\overline K})$ est quasi unipotente ssi son action sur $H^i_!(U_{\overline K})$ l'est : caractère \uline{birationnel} de la question.

Si donc on peut trouver un \deleted{\unsure{modèle}} \inserted{``complété'' propre et lisse $\widehat U_K$} \deleted{\unsure{lisse}} de \deleted{\unsure{$X_K$}} $U_K$, alors on est ramené au cas d'un schéma \uline{lisse} et \uline{propre}. Si on admet la résolution des singularités \deleted{\unsure{par}} \inserted{(\unsure{par voie forte})}, sur la clôture \deleted{\unsure{alg.}} \inserted{parfaite} de $K$, pour les schémas de t.f. de dim $\leq n$, on aura gagné \inserted{(prenant $U$ affine)}, à condition que l'hypothèse de récurrence reste vérifiée après toute extension \inserted{radicale} finie \deleted{\unsure{ext}} $K'$ de $K$ (en prenant la clôture normale $V'$ de $V$ dans $K'$).

\inserted{\uline{Remarque} 9.3.} Donc en vertu de \deleted{\unsure{n°}} \unsure{7.1|9.1}, ça marche aussi pour $H^i(X_{\overline K},\unsure{\mathbb Q_\ell|Q_\ell})$, si $X$ lisse \uline{séparé}. Le cas non séparé s'y ramène, par suite spectrale de Leray d'un recouvrement fini de $X$ par des ouverts affines.

\deleted{9.3. Cas de $H^i_!(X_{\overline K},\mathbb Q_\ell)$ pour $X$ \inserted{séparé de t.f.} quelconque. \deleted{\unsure{Par ailleurs},} \deleted{On peut supposer $X$ réduit, et si} $U$ un ouvert dense, on \deleted{\unsure{peut}} \unsure{est} \unsure{comme} \deleted{\unsure{cela|alors}} \unsure{précédemment} \unsure{quitte pour|quitte à} remplacer $X$ par $U$. On peut trouver $U$ tel que \deleted{\unsure{changer $U$}} sur une \deleted{\unsure{une}} extension \inserted{finie} radicale $K'$ convenable de $K$, on ait $(\deleted{\unsure{X}}U_{K'})_{\mathrm{red}}$ \uline{lisse} sur $K'$. On est ramené au cas traité dans 9.2.}

\deleted{9.3. \uline{Cas de $H^i(X_{\overline K},\mathbb Q_\ell)$ pour $X$ de t.f. quelconque.}}
